% -------------------------------------------------------
%  Minimal Persian-command compatibility layer
% -------------------------------------------------------
%
% The thesis originally depended on XePersian's legacy `localise=on`
% option. Only the aliases used by this project are defined here, allowing
% XePersian to be loaded without that version-sensitive option.

% Persian compound command names in this project contain U+200C (ZWNJ).
% Treat it as a letter, matching the behaviour of XePersian's old localise
% layer, so names such as \بدون‌تورفتگی remain a single control sequence.
\catcode"200C=11\relax

% Core LaTeX commands used in the thesis sources.
\let\شروع\begin
\let\پایان\end
\let\فصل\chapter
\let\قسمت\section
\let\زیرقسمت\subsection
\let\مرجع\cite
\let\رجوع\ref
\let\برچسب\label
\let\شرح\caption
\let\فقره\item
\let\کوچک\small
\let\بدون‌تورفتگی\noindent
\let\خط‌پر\hline

% Environment names used through \شروع{...} and \پایان{...}.
\newenvironment{شکل}{\begin{figure}}{\end{figure}}
\newenvironment{لوح}{\begin{table}}{\end{table}}
\newenvironment{جدول}{\begin{tabular}}{\end{tabular}}
\newenvironment{فقرات}{\begin{itemize}}{\end{itemize}}
\newenvironment{شمارش}{\begin{enumerate}}{\end{enumerate}}
\newenvironment{وسط‌چین}{\begin{center}}{\end{center}}
